\begin{proposition}[Operator-norm convergence of the metric operator]
\label{prop:operator-norm-convergence}
Under the closed embedded-manifold and smooth positive-density
assumptions of the main theorem, write
\[
 A=-\Delta_M-\nabla\log\rho\cdot\nabla,
 \qquad G=A^{-1},\qquad
 B_\varepsilon=
   \frac{\varepsilon}{2}(I-P_\varepsilon^2)^{-1}P_\varepsilon,
\]
where inverses are taken on $\mathcal H=L^2_0(\mu)$.
Then
\[
 \|B_\varepsilon-G\|_{\mathcal H\to\mathcal H}
 \longrightarrow0.
\]
This assertion supplies no operator-norm rate or higher-order
operator-norm expansion.
\end{proposition}

\begin{proof}
We first strengthen the spectral-gap comparison to a quadratic-form
bound by the weighted heat semigroup.  Put $\sigma=\sqrt\rho$.
Multiplication by $\sigma$ is a unitary map from $L^2(\mu)$ to
$L^2(\mathrm{vol}_M)$, under which
\[
 \sigma A\sigma^{-1}=-\Delta_M+V,
 \qquad V=\frac{\Delta_M\sigma}{\sigma}.
\]
The smooth potential $V$ is bounded.  If $h_t$ is the ordinary heat
kernel and $h_t^V$ is the kernel of $e^{-t(-\Delta_M+V)}$, the
maximum principle gives
\[
 0\le h_t^V(x,y)\le e^{t\|V_-\|_\infty}h_t(x,y).
\]
Consequently the $\mu$-kernel $q_t$ of $H_t=e^{-tA}$ satisfies
\[
 q_t(x,y)=\frac{h_t^V(x,y)}{\sigma(x)\sigma(y)}
 \le C t^{-m/2}
       \exp\!\left(-\frac{d_M(x,y)^2}{C_0t}\right)
\]
for sufficiently small $t>0$.  Here we used the usual small-time
Gaussian upper bound on a compact closed manifold and the positive
lower bound on $\sigma$.

The already proved uniform scaling asymptotic gives
$0<v_*\le v_\varepsilon(x)\le v^*<\infty$, where
$v_\varepsilon=(\pi\varepsilon)^{m/4}u_\varepsilon$.
The $\mu$-kernel of $P_\varepsilon$ is
\[
 k_\varepsilon(x,y)=
 v_\varepsilon(x)v_\varepsilon(y)
 (\pi\varepsilon)^{-m/2}e^{-|x-y|^2/\varepsilon}.
\]
Choose $t=\alpha\varepsilon$, with $\alpha=1/C_0$.
Because $|x-y|\le d_M(x,y)$, the preceding bounds imply
$k_\varepsilon(x,y)\ge c\,q_{\alpha\varepsilon}(x,y)$ for some
$c>0$ independent of sufficiently small $\varepsilon$.
Both kernels are symmetric Markov kernels with respect to $\mu$.
Their Dirichlet-form representations therefore give, on
$\mathcal H$,
\[
 I-P_\varepsilon\ \ge\ c(I-e^{-\alpha\varepsilon A}).
\]
Both operators are strictly positive and invertible on this
mean-zero space.  Inversion reverses the positive-operator order;
no commutativity between these two operators is required.  Moreover,
the functional calculus of $P_\varepsilon$ gives
\[
 0\le B_\varepsilon
 =\frac{\varepsilon}{4}
  \bigl[(I-P_\varepsilon)^{-1}-(I+P_\varepsilon)^{-1}\bigr]
 \le\frac{\varepsilon}{4c}
       (I-e^{-\alpha\varepsilon A})^{-1}.
\]
For $z>0$ one has $(1-e^{-z})^{-1}\le1+z^{-1}$.
Applying this scalar inequality to the spectral calculus of $A$
proves the compact-dominance estimate
\begin{equation}
\label{eq:compact-dominance}
 0\le B_\varepsilon
 \le\frac1{4c\alpha}G+\frac{\varepsilon}{4c}I
 \le C(G+\varepsilon I).
\end{equation}

Let $0<\lambda_1\le\lambda_2\le\cdots\to\infty$ be the
eigenvalues of $A$ on $\mathcal H$, counted with multiplicity.
Let $E_J$ project onto the first $J$ eigenfunctions and set
$F_J=I-E_J$.  These eigenfunctions are smooth.  The previously
proved convergence on smooth directions, together with finite
dimensionality, gives
\[
 \|(B_\varepsilon-G)E_J\|\longrightarrow0
 \quad\text{for each fixed }J.
\]
Writing $D_\varepsilon=B_\varepsilon-G$, self-adjointness and
\[
 D_\varepsilon
 =D_\varepsilon E_J+E_JD_\varepsilon F_J
   +F_JD_\varepsilon F_J
\]
yield
\[
 \|D_\varepsilon\|
 \le2\|D_\varepsilon E_J\|
      +\|F_JD_\varepsilon F_J\|.
\]
Indeed,
$\|E_JD_\varepsilon F_J\|
=\|F_JD_\varepsilon E_J\|\le\|D_\varepsilon E_J\|$.
By \eqref{eq:compact-dominance},
\[
 \|F_JD_\varepsilon F_J\|
 \le\|F_JB_\varepsilon F_J\|+\|F_JGF_J\|
 \le C\varepsilon+\frac{C+1}{\lambda_{J+1}}.
\]
First let $\varepsilon\downarrow0$ for fixed $J$, and then let
$J\to\infty$.  The resulting upper bound tends to zero, proving
the proposition.
\end{proof}

\paragraph{Scope.}
The compact-dominance and finite-dimensional projection argument is
a general operator-theoretic device, not a separate novelty claim.
The proposition strengthens strong convergence of $B_\varepsilon$
to norm convergence. The subsequent uniform gradient lemma in the final manuscript also transfers norm convergence to the positive correction. The smooth-direction asymptotic expansion does not become an operator-norm expansion.  In particular it is
consistent with the circle's order-$\varepsilon$ lower bound.
